\chapter{Nodes, Items und Craften}
\label{ch:items/nodes-items-crafting}

Neue Nodes und Craftitems\index{craftitem@Craftitem} zu registrieren und
Craftingrezepte\index{craftingrezept@Craftingrezept} anzulegen, gehört
zu den Grundlagen vieler Mods.

\section{Was sind Nodes und Items?}\label{was-sind-nodes-und-items}

Nodes, Craftitems und Werkzeuge sind alles Items. Ein Item ist etwas,
das sich in einem Inventar\index{inventar@Inventar} befinden kann --
selbst wenn das im normalen Spielverlauf nicht möglich ist.

Ein Node ist ein Item, das platziert oder in der Welt gefunden werden
kann. Jede Position in der Welt ist mit genau einem Node belegt --
scheinbar leere Positionen sind normalerweise Luft-Nodes.

Ein Craftitem kann nicht platziert werden. Es kommt nur in Inventaren
oder als fallengelassenes Item in der Welt vor.

Ein Werkzeug ist wie ein Craftitem, kann sich aber abnutzen. Während du
das Werkzeug benutzt, sinkt der Abnutzungsbalken, bis das Werkzeug
zerbricht. Werkzeuge lassen sich außerdem nie stapeln. Vermutlich werden
Craftitems und Werkzeuge künftig zu einer einzigen Item-Art
verschmelzen, da die Unterscheidung zwischen ihnen recht künstlich ist.

\section{Items registrieren}\label{items-registrieren}

Eine Item-Definition besteht aus einem \emph{Itemnamen} und einer
\emph{Definitionstabelle}. Die Definitionstabelle enthält Attribute, die
das Verhalten des Items bestimmen.

\begin{Shaded}
\begin{Highlighting}[]
\VariableTok{core}\OperatorTok{.}\NormalTok{register\_craftitem}\OperatorTok{(}\StringTok{"modname:itemname"}\OperatorTok{,} \OperatorTok{\{}
    \VariableTok{description} \OperatorTok{=} \StringTok{"Mein besonderes Item"}\OperatorTok{,}
    \VariableTok{inventory\_image} \OperatorTok{=} \StringTok{"modname\_itemname.png"}\OperatorTok{,}
\OperatorTok{\})}
\end{Highlighting}
\end{Shaded}

\subsection{Itemnamen}\label{itemnamen}

Jedes Item hat einen Itemnamen, über den es angesprochen wird. Er hat
folgendes Format:

\begin{Shaded}
\begin{Highlighting}[]
\NormalTok{modname:itemname}
\end{Highlighting}
\end{Shaded}

\texttt{modname} ist der Name des Mods, in dem das Item registriert
wird, \texttt{itemname} der Name des Items selbst. Der Itemname sollte
beschreiben, was das Item ist, und darf nicht bereits registriert sein.

Sowohl \texttt{modname} als auch \texttt{itemname} sollten nur
Kleinbuchstaben, Ziffern und Unterstriche enthalten.

\subsection{Item-Aliase}\label{item-aliase}

Items können auch \emph{Aliase\index{alias@Alias}} haben, die auf ihren
Namen verweisen. Ein \emph{Alias} ist ein Pseudo-Itemname: Die Engine
behandelt jedes Vorkommen des Alias so, als stünde dort der eigentliche
Itemname. Dafür gibt es zwei häufige Anwendungsfälle:

\begin{itemize}
\tightlist
\item
  Entfernte Items umbenennen. Wird ein Item ohne Korrekturcode aus einem
  Mod entfernt, können unbekannte Nodes in der Welt und in Inventaren
  zurückbleiben.
\item
  Abkürzungen anlegen. \texttt{/\allowbreak{}giveme dirt} ist einfacher
  als \texttt{/\allowbreak{}giveme default:dirt}.
\end{itemize}

Einen Alias zu registrieren, ist einfach. Die Reihenfolge der Argumente
merkst du dir am besten als \texttt{von → nach}, wobei \emph{von} der
Alias und \emph{nach} das Ziel ist.

\begin{Shaded}
\begin{Highlighting}[]
\VariableTok{core}\OperatorTok{.}\NormalTok{register\_alias}\OperatorTok{(}\StringTok{"dirt"}\OperatorTok{,} \StringTok{"default:dirt"}\OperatorTok{)}
\end{Highlighting}
\end{Shaded}

Mods müssen Aliase selbst auflösen, bevor sie direkt mit Itemnamen
arbeiten, denn die Engine tut das nicht. Das geht aber ganz einfach:

\begin{Shaded}
\begin{Highlighting}[]
\VariableTok{itemname} \OperatorTok{=} \VariableTok{core}\OperatorTok{.}\VariableTok{registered\_aliases}\OperatorTok{[}\VariableTok{itemname}\OperatorTok{]} \KeywordTok{or} \VariableTok{itemname}
\end{Highlighting}
\end{Shaded}

\subsection{Texturen}\label{texturen}

Texturen\index{textur@Textur} gehören in den Ordner
\texttt{textures/\allowbreak{}} und tragen Namen im Format
\texttt{modname\_\allowbreak{}itemname.\allowbreak{}png}. JPEG-Texturen
werden unterstützt, können aber keine Transparenz und sehen bei
niedrigen Auflösungen meist schlecht aus. Meist ist PNG die bessere
Wahl.

Texturen in Luanti sind in der Regel 16 × 16 Pixel groß. Sie können jede
Auflösung haben, empfohlen sind aber Zweierpotenzen, etwa 16, 32, 64
oder 128. Andere Auflösungen werden auf älteren Geräten, besonders auf
Smartphones, möglicherweise nicht richtig unterstützt und verschlechtern
dann die Leistung.

\section{Einen einfachen Node
registrieren}\label{einen-einfachen-node-registrieren}

Nodes werden ähnlich wie Items registriert, nur mit einer anderen
Funktion:

\begin{Shaded}
\begin{Highlighting}[]
\VariableTok{core}\OperatorTok{.}\NormalTok{register\_node}\OperatorTok{(}\StringTok{"mymod:diamond"}\OperatorTok{,} \OperatorTok{\{}
    \VariableTok{description} \OperatorTok{=} \StringTok{"Alien{-}Diamant"}\OperatorTok{,}
    \VariableTok{tiles} \OperatorTok{=} \OperatorTok{\{} \StringTok{"mymod\_diamond.png"} \OperatorTok{\},}
    \VariableTok{is\_ground\_content} \OperatorTok{=} \KeywordTok{true}\OperatorTok{,}
    \VariableTok{groups} \OperatorTok{=} \OperatorTok{\{} \VariableTok{cracky} \OperatorTok{=} \DecValTok{3}\OperatorTok{,} \VariableTok{stone} \OperatorTok{=} \DecValTok{1} \OperatorTok{\},}
\OperatorTok{\})}
\end{Highlighting}
\end{Shaded}

Node-Definitionen können jede Eigenschaft einer Item-Definition
enthalten und haben zusätzlich eigene, node-spezifische Eigenschaften.

Die Eigenschaft \texttt{tiles} ist eine Tabelle mit den Texturnamen, die
der Node verwendet. Gibt es nur eine Textur, wird sie auf allen Seiten
verwendet. Für eine eigene Textur je Seite gibst du die Namen von sechs
Texturen in dieser Reihenfolge an:

\begin{Shaded}
\begin{Highlighting}[]
\NormalTok{oben (+Y), unten ({-}Y), rechts (+X), links ({-}X),}
\NormalTok{hinten (+Z), vorne ({-}Z)}
\end{Highlighting}
\end{Shaded}

Denk daran, dass +Y in Luanti nach oben zeigt, wie in den meisten
3D-Computerspielen üblich.

\begin{Shaded}
\begin{Highlighting}[]
\VariableTok{core}\OperatorTok{.}\NormalTok{register\_node}\OperatorTok{(}\StringTok{"mymod:diamond"}\OperatorTok{,} \OperatorTok{\{}
    \VariableTok{description} \OperatorTok{=} \StringTok{"Alien{-}Diamant"}\OperatorTok{,}
    \VariableTok{tiles} \OperatorTok{=} \OperatorTok{\{}
        \StringTok{"mymod\_diamond\_up.png"}\OperatorTok{,} \CommentTok{{-}{-} y+}
        \StringTok{"mymod\_diamond\_down.png"}\OperatorTok{,} \CommentTok{{-}{-} y{-}}
        \StringTok{"mymod\_diamond\_right.png"}\OperatorTok{,} \CommentTok{{-}{-} x+}
        \StringTok{"mymod\_diamond\_left.png"}\OperatorTok{,} \CommentTok{{-}{-} x{-}}
        \StringTok{"mymod\_diamond\_back.png"}\OperatorTok{,} \CommentTok{{-}{-} z+}
        \StringTok{"mymod\_diamond\_front.png"}\OperatorTok{,} \CommentTok{{-}{-} z{-}}
    \OperatorTok{\},}
    \VariableTok{is\_ground\_content} \OperatorTok{=} \KeywordTok{true}\OperatorTok{,}
    \VariableTok{groups} \OperatorTok{=} \OperatorTok{\{} \VariableTok{cracky} \OperatorTok{=} \DecValTok{3} \OperatorTok{\},}
    \VariableTok{drop} \OperatorTok{=} \StringTok{"mymod:diamond\_fragments"}\OperatorTok{,}
    \CommentTok{{-}{-} \^{} statt des Diamanten Fragmente fallen lassen}
\OperatorTok{\})}
\end{Highlighting}
\end{Shaded}

Das Attribut \texttt{is\_\allowbreak{}ground\_\allowbreak{}content}
erlaubt es, Höhlen durch diesen Node zu graben. Das ist wichtig für
jeden Node, der bei der
Kartengenerierung\index{kartengenerierung@Kartengenerierung} unter der
Erde platziert werden kann. Höhlen werden aus der Welt
herausgeschnitten, nachdem alle anderen Nodes eines Gebiets generiert
wurden.

\section{Crafting}\label{crafting}

Es gibt mehrere Arten von Craftingrezepten; welche gemeint ist, gibt die
Eigenschaft \texttt{type} an.

\begin{itemize}
\tightlist
\item
  \texttt{shaped} -- Die Zutaten müssen an der richtigen Stelle liegen.
\item
  \texttt{shapeless} -- Es ist egal, wo die Zutaten liegen; nur die
  Menge muss stimmen.
\item
  \texttt{cooking} -- Rezepte für den Ofen.
\item
  \texttt{fuel} -- Legt fest, welche Items in Öfen verbrannt werden
  können.
\item
  \texttt{tool\_\allowbreak{}repair} -- Legt fest, welche Items sich als
  Werkzeug reparieren lassen.
\end{itemize}

Craftingrezepte sind keine Items und haben daher keinen Itemnamen, der
sie eindeutig kennzeichnet.

\subsection{Shaped}\label{shaped}

Bei Shaped-Rezepten müssen die Zutaten in der richtigen Form, also im
richtigen Muster, angeordnet sein. Im folgenden Beispiel müssen die
Fragmente wie ein Stuhl angeordnet sein, damit das Rezept funktioniert.
Die lokale Variable \texttt{frag} hält das Rezept übersichtlich.

\begin{Shaded}
\begin{Highlighting}[]
\KeywordTok{local} \VariableTok{frag} \OperatorTok{=} \StringTok{"mymod:diamond\_fragments"}

\VariableTok{core}\OperatorTok{.}\NormalTok{register\_craft}\OperatorTok{(\{}
    \FunctionTok{type} \OperatorTok{=} \StringTok{"shaped"}\OperatorTok{,}
    \VariableTok{output} \OperatorTok{=} \StringTok{"mymod:diamond\_chair 99"}\OperatorTok{,}
    \VariableTok{recipe} \OperatorTok{=} \OperatorTok{\{}
        \OperatorTok{\{} \VariableTok{frag}\OperatorTok{,} \StringTok{""}\OperatorTok{,}   \StringTok{""} \OperatorTok{\},}
        \OperatorTok{\{} \VariableTok{frag}\OperatorTok{,} \VariableTok{frag}\OperatorTok{,} \StringTok{""} \OperatorTok{\},}
        \OperatorTok{\{} \VariableTok{frag}\OperatorTok{,} \VariableTok{frag}\OperatorTok{,} \StringTok{""} \OperatorTok{\},}
    \OperatorTok{\},}
\OperatorTok{\})}
\end{Highlighting}
\end{Shaded}

Beachte die leere Spalte auf der rechten Seite. Sie bedeutet, dass
rechts von der Form eine leere Spalte vorhanden sein \emph{muss}, sonst
funktioniert das Rezept nicht. Soll diese leere Spalte nicht nötig sein,
lässt du die leeren Strings einfach weg:

\begin{Shaded}
\begin{Highlighting}[]
\VariableTok{core}\OperatorTok{.}\NormalTok{register\_craft}\OperatorTok{(\{}
    \VariableTok{output} \OperatorTok{=} \StringTok{"mymod:diamond\_chair 99"}\OperatorTok{,}
    \VariableTok{recipe} \OperatorTok{=} \OperatorTok{\{}
        \OperatorTok{\{} \VariableTok{frag}\OperatorTok{,} \StringTok{""}   \OperatorTok{\},}
        \OperatorTok{\{} \VariableTok{frag}\OperatorTok{,} \VariableTok{frag} \OperatorTok{\},}
        \OperatorTok{\{} \VariableTok{frag}\OperatorTok{,} \VariableTok{frag} \OperatorTok{\},}
    \OperatorTok{\},}
\OperatorTok{\})}
\end{Highlighting}
\end{Shaded}

Das Feld \texttt{type} wird für Shaped-Rezepte eigentlich nicht
gebraucht, da \texttt{shaped} der Standardtyp ist.

\subsection{Shapeless}\label{shapeless}

Shapeless-Rezepte verwendest du, wenn es egal ist, wo die Zutaten liegen
-- Hauptsache, sie sind da.

\begin{Shaded}
\begin{Highlighting}[]
\VariableTok{core}\OperatorTok{.}\NormalTok{register\_craft}\OperatorTok{(\{}
    \FunctionTok{type} \OperatorTok{=} \StringTok{"shapeless"}\OperatorTok{,}
    \VariableTok{output} \OperatorTok{=} \StringTok{"mymod:diamond 3"}\OperatorTok{,}
    \VariableTok{recipe} \OperatorTok{=} \OperatorTok{\{}
        \StringTok{"mymod:diamond\_fragments"}\OperatorTok{,}
        \StringTok{"mymod:diamond\_fragments"}\OperatorTok{,}
        \StringTok{"mymod:diamond\_fragments"}\OperatorTok{,}
    \OperatorTok{\},}
\OperatorTok{\})}
\end{Highlighting}
\end{Shaded}

\subsection{Cooking und Fuel}\label{cooking-und-fuel}

Rezepte vom Typ \texttt{cooking} werden nicht im Craftinggitter
hergestellt, sondern in Öfen oder anderen Kochgeräten gekocht, die es in
Mods geben kann.

\begin{Shaded}
\begin{Highlighting}[]
\VariableTok{core}\OperatorTok{.}\NormalTok{register\_craft}\OperatorTok{(\{}
    \FunctionTok{type} \OperatorTok{=} \StringTok{"cooking"}\OperatorTok{,}
    \VariableTok{output} \OperatorTok{=} \StringTok{"mymod:diamond\_fragments"}\OperatorTok{,}
    \VariableTok{recipe} \OperatorTok{=} \StringTok{"default:coalblock"}\OperatorTok{,}
    \VariableTok{cooktime} \OperatorTok{=} \DecValTok{10}\OperatorTok{,}
\OperatorTok{\})}
\end{Highlighting}
\end{Shaded}

Der einzige echte Unterschied im Code ist, dass \texttt{recipe} nur ein
einzelnes Item ist statt einer Tabelle (in geschweiften Klammern).
Optional gibt es den Parameter \texttt{cooktime}, der festlegt, wie
lange das Item zum Kochen braucht. Ohne Angabe sind es 3 Sekunden.

Das Rezept oben funktioniert, wenn der Kohleblock im Eingabefeld liegt
und darunter irgendein Brennstoff. Nach 10 Sekunden entstehen
Diamantfragmente!

Der Typ \texttt{fuel} ergänzt den Typ \texttt{cooking}: Er legt fest,
was in Öfen und anderen Kochgeräten aus Mods verbrannt werden kann.

\begin{Shaded}
\begin{Highlighting}[]
\VariableTok{core}\OperatorTok{.}\NormalTok{register\_craft}\OperatorTok{(\{}
    \FunctionTok{type} \OperatorTok{=} \StringTok{"fuel"}\OperatorTok{,}
    \VariableTok{recipe} \OperatorTok{=} \StringTok{"mymod:diamond"}\OperatorTok{,}
    \VariableTok{burntime} \OperatorTok{=} \DecValTok{300}\OperatorTok{,}
\OperatorTok{\})}
\end{Highlighting}
\end{Shaded}

Fuel-Rezepte haben kein Ergebnis wie andere Rezepte, dafür aber eine
Brenndauer in Sekunden, die angibt, wie lange sie als Brennstoff
reichen. Der Diamant brennt also 300 Sekunden lang!

\section{Gruppen}\label{gruppen}

Items können Mitglied in vielen Gruppen\index{gruppe@Gruppe} sein, und
Gruppen können viele Mitglieder haben. Gruppen werden mit der
Eigenschaft \texttt{groups} in der Definitionstabelle festgelegt und
haben jeweils einen zugehörigen Wert.

\begin{Shaded}
\begin{Highlighting}[]
\VariableTok{groups} \OperatorTok{=} \OperatorTok{\{} \VariableTok{cracky} \OperatorTok{=} \DecValTok{3}\OperatorTok{,} \VariableTok{wood} \OperatorTok{=} \DecValTok{1} \OperatorTok{\}}
\end{Highlighting}
\end{Shaded}

Gruppen haben mehrere Zwecke. Erstens beschreiben sie Eigenschaften wie
Abbautyp\index{abbautyp@Abbautyp} und Brennbarkeit. Zweitens kann in
einem Craftingrezept statt eines Itemnamens eine Gruppe stehen; dann
passt jedes Item dieser Gruppe.

\begin{Shaded}
\begin{Highlighting}[]
\VariableTok{core}\OperatorTok{.}\NormalTok{register\_craft}\OperatorTok{(\{}
    \FunctionTok{type} \OperatorTok{=} \StringTok{"shapeless"}\OperatorTok{,}
    \VariableTok{output} \OperatorTok{=} \StringTok{"mymod:diamond\_thing 3"}\OperatorTok{,}
    \VariableTok{recipe} \OperatorTok{=} \OperatorTok{\{} \StringTok{"group:wood"}\OperatorTok{,} \StringTok{"mymod:diamond"} \OperatorTok{\},}
\OperatorTok{\})}
\end{Highlighting}
\end{Shaded}

\section{Werkzeuge, Fähigkeiten und
Abbautypen}\label{werkzeuge-fuxe4higkeiten-und-abbautypen}

Abbautypen sind Gruppen, die festlegen, wie widerstandsfähig ein Node
beim Abbauen mit verschiedenen Werkzeugen ist. Ein höherer Wert in einer
Abbautyp-Gruppe bedeutet, dass sich der Node leichter und schneller
abbauen lässt. Mehrere Abbautypen lassen sich kombinieren, damit
verschiedene Werkzeugarten effizient eingesetzt werden können. Ein Node
ohne Abbautyp kann mit keinem Werkzeug abgebaut werden.

\begin{longtable}[]{@{}
  >{\raggedright\arraybackslash}p{(\linewidth - 4\tabcolsep) * \real{0.2051}}
  >{\raggedright\arraybackslash}p{(\linewidth - 4\tabcolsep) * \real{0.4359}}
  >{\raggedright\arraybackslash}p{(\linewidth - 4\tabcolsep) * \real{0.3590}}@{}}
\toprule\noalign{}
\begin{minipage}[b]{\linewidth}\raggedright
Gruppe
\end{minipage} & \begin{minipage}[b]{\linewidth}\raggedright
Bestes Werkzeug
\end{minipage} & \begin{minipage}[b]{\linewidth}\raggedright
Beschreibung
\end{minipage} \\
\midrule\noalign{}
\endhead
\bottomrule\noalign{}
\endlastfoot
\texttt{crumbly} & Schaufel & Erde, Sand \\
\texttt{cracky} & Spitzhacke & Zähes (aber sprödes) Material wie
Stein \\
\texttt{snappy} & \emph{beliebig} & Lässt sich mit feinen Werkzeugen
schneiden, z.~B. Blätter, kleine Pflanzen, Draht, Blech \\
\texttt{choppy} & Axt & Lässt sich mit scharfer Kraft schneiden, z.~B.
Bäume, Holzbretter \\
\texttt{fleshy} & Schwert & Lebewesen wie Tiere und der Spieler; kann
Bluteffekte beim Treffen nahelegen \\
\texttt{explody} & ? & Besonders anfällig für Explosionen \\
\texttt{oddly\_\allowbreak{}breakable\_\allowbreak{}by\_\allowbreak{}hand}
& \emph{beliebig} & Fackeln und Ähnliches -- sehr schnell abzubauen \\
\end{longtable}

Jedes Werkzeug hat eine Werkzeugfähigkeit (\emph{tool capability}). Sie
umfasst eine Liste der unterstützten Abbautypen und für jeden Typ
zugehörige Eigenschaften wie Abbauzeiten und
Abnutzung.\index{abnutzung@Abnutzung} Werkzeuge können außerdem je Typ
eine maximale Härte haben; so lässt sich verhindern, dass schwächere
Werkzeuge härtere Nodes abbauen. Sehr häufig enthalten Werkzeuge alle
Abbautypen in ihren Fähigkeiten, wobei die weniger geeigneten sehr
ineffiziente Werte haben. Hat das Item, das ein Spieler gerade in der
Hand hält, keine eigene Werkzeugfähigkeit, wird stattdessen die
Fähigkeit der Hand verwendet.

\begin{Shaded}
\begin{Highlighting}[]
\VariableTok{core}\OperatorTok{.}\NormalTok{register\_tool}\OperatorTok{(}\StringTok{"mymod:tool"}\OperatorTok{,} \OperatorTok{\{}
    \VariableTok{description} \OperatorTok{=} \StringTok{"Mein Werkzeug"}\OperatorTok{,}
    \VariableTok{inventory\_image} \OperatorTok{=} \StringTok{"mymod\_tool.png"}\OperatorTok{,}
    \VariableTok{tool\_capabilities} \OperatorTok{=} \OperatorTok{\{}
        \VariableTok{full\_punch\_interval} \OperatorTok{=} \DecValTok{1.5}\OperatorTok{,}
        \VariableTok{max\_drop\_level} \OperatorTok{=} \DecValTok{1}\OperatorTok{,}
        \VariableTok{groupcaps} \OperatorTok{=} \OperatorTok{\{}
            \VariableTok{crumbly} \OperatorTok{=} \OperatorTok{\{}
                \VariableTok{maxlevel} \OperatorTok{=} \DecValTok{2}\OperatorTok{,}
                \VariableTok{uses} \OperatorTok{=} \DecValTok{20}\OperatorTok{,}
                \VariableTok{times} \OperatorTok{=} \OperatorTok{\{} \OperatorTok{[}\DecValTok{1}\OperatorTok{]} \OperatorTok{=} \DecValTok{1.60}\OperatorTok{,} \OperatorTok{[}\DecValTok{2}\OperatorTok{]} \OperatorTok{=} \DecValTok{1.20}\OperatorTok{,} \OperatorTok{[}\DecValTok{3}\OperatorTok{]} \OperatorTok{=} \DecValTok{0.80} \OperatorTok{\},}
            \OperatorTok{\},}
        \OperatorTok{\},}
        \VariableTok{damage\_groups} \OperatorTok{=} \OperatorTok{\{} \VariableTok{fleshy} \OperatorTok{=} \DecValTok{2} \OperatorTok{\},}
    \OperatorTok{\},}
\OperatorTok{\})}
\end{Highlighting}
\end{Shaded}

\texttt{groupcaps} ist die Liste der Abbautypen, mit denen das Werkzeug
Nodes abbauen kann. Schadensgruppen\index{schadensgruppe@Schadensgruppe}
(\emph{damage groups}) steuern, wie Werkzeuge Objekten Schaden zufügen;
das wird später im Kapitel „Objekte, Spieler und
Entities``\index{entity@Entity} behandelt.
